\section{A Kummer--torsion specialization theorem}
\label{sec:abstract-specialization}

For a number field $B$, write
\[
 B^c=B(\mu_\infty).
\]
Thus $\Q^c=\Q^{\mathrm{ab}}$ by Kronecker--Weber.  We first record a
one-dimensional form of cyclotomic Hilbert irreducibility that will be
used twice below.

Let $X/\Q$ be a geometrically integral normal curve and let
$t,y\in\Q(X)$ satisfy
\[
 \Q(X)=\Q(t,y),\qquad
 N=[\Q(X):\Q(y)].
\]
Regard $y$ as a finite cover of $\mathbf G_m$ after deleting its zeros,
poles, and finitely many other points.  Let
\[
 A(T,Y)\in\Q(Y)[T]
\]
be the monic minimal polynomial of $t$ over $\Q(Y)$.  There is a finite
Galois-stable set $\Sigma_{\rm alg}\subset\mu_\infty$ such that, for
$\theta\notin\Sigma_{\rm alg}$, specialization at $Y=\theta$ is
defined, $A(T,\theta)$ is separable of degree $N$, and its roots are
the $t$-coordinates of the points of $y^{-1}(\theta)$.  Indeed, this
follows by shrinking the target so that the cover is finite \'etale and
the primitive-element plane model is isomorphic to it.

We say that $t$ separates the torsion fibers of $y$ away from a finite
set if there is a finite Galois-stable
$\Sigma_{\rm sep}\subset\mu_\infty$ such that
\[
 P\longmapsto t(P)
 \numberthis\label{eq:torsion-separation}
\]
is injective on
\[
 \bigcup_{\theta\in\mu_\infty\setminus\Sigma_{\rm sep}}
 y^{-1}(\theta)(\overline{\Q}).
\]
This single condition has two logically distinct consequences: fibers
over different torsion points have disjoint sets of $t$-coordinates,
and the torsion label can be recovered from any one such coordinate.

Recall that a cover $y:X\to\mathbf G_m$ satisfies the pullback
condition \emph{(PB)} if its pullback by $u\mapsto u^n$ is geometrically
irreducible for every $n\ge1$ \cite[Definition and Proposition
2.1]{Zannier2010}.  Equivalently,
\[
 U^n-y\quad\hbox{is irreducible in }\overline{\Q}(X)[U]
 \quad(n\ge1).
 \numberthis\label{eq:pb-binomial}
\]

\begin{theorem}[Kummer--torsion specialization]
\label{thm:kummer-torsion}
Assume that $y$ satisfies \emph{(PB)} and that $t$ separates its torsion
fibers away from a finite set.  Fix $r\ge1$ and put $x=y^r$.  For
$k\ge2$, set
\[
 S_r(k)=\{\theta\in\mu_\infty:\ord(\theta^r)=k\}
\]
and, whenever all the specializations are good, put
\[
 \Psi_{r,k}(T)=\prod_{\theta\in S_r(k)}A(T,\theta).
 \numberthis\label{eq:abstract-Psi}
\]
Then the following assertions hold for every sufficiently large $k$.

\begin{enumerate}[label=\textup{(\roman*)}]
\item Over $\Q^{\mathrm{ab}}$, the polynomial $\Psi_{r,k}$ has exactly
      $r\varphi(k)$ distinct irreducible factors, namely the
      $A(T,\theta)$, and every factor has degree $N$.
\item The possible orders of elements of $S_r(k)$ are precisely
      \[
       n=ka,\qquad a\mid r,\qquad (k,r/a)=1.
       \numberthis\label{eq:admissible-orders}
      \]
      For every such $a$, the polynomial
      \[
       F_{k,a}(T)=
       \prod_{\ord(\theta)=ka}A(T,\theta)
       \numberthis\label{eq:order-block}
      \]
      belongs to $\Q[T]$, is irreducible there, and has degree
      $N\varphi(ka)$.  The product of these $F_{k,a}$ is the complete
      factorization of $\Psi_{r,k}$ over $\Q$.
\item If $t_0$ is a root of $A(T,\theta)$, then
      \[
       \theta\in\Q(t_0),\qquad
       \Q(t_0)\cap\Q^{\mathrm{ab}}=\Q(\theta).
       \numberthis\label{eq:abstract-intersection}
      \]
\end{enumerate}
\end{theorem}

\begin{proof}
By Zannier's cyclotomic Hilbert irreducibility theorem
\cite[Theorem 2.1]{Zannier2010}, condition (PB) gives a finite set
$\Sigma_{\rm Zan}\subset\mu_\infty$ outside which the fiber of $y$ is
irreducible over $\Q^c=\Q^{\mathrm{ab}}$ and has degree $N$.  Enlarge
this set by $\Sigma_{\rm alg}$ and $\Sigma_{\rm sep}$.  Since
$|S_r(k)|=r\varphi(k)$ and every $\theta\in S_r(k)$ has order at least
$k$, no element of this fixed finite set occurs once $k$ is large.
This proves the irreducibility and degree assertion in (i); separation
shows that the factors have disjoint root sets.

Let $P\in y^{-1}(\theta)$ and write $t_0=t(P)$.  If
$\sigma\in\Gal(\overline\Q/\Q(t_0))$, then
\[
 t(\sigma P)=t_0,\qquad y(\sigma P)=\sigma\theta.
\]
Torsion separation forces $\sigma P=P$, and hence
$\sigma\theta=\theta$.  Thus $\theta\in\Q(t_0)$.  Irreducibility over
$\Q^{\mathrm{ab}}$ gives
\[
 [\Q^{\mathrm{ab}}(t_0):\Q^{\mathrm{ab}}]=N.
\]
It also implies that $A(T,\theta)$ is irreducible over $\Q(\theta)$,
so $[\Q(t_0):\Q(\theta)]=N$.  If
$D=\Q(t_0)\cap\Q^{\mathrm{ab}}$, the compositum degree formula gives
$[\Q(t_0):D]=N$.  Since $\Q(\theta)\subseteq D$, the two degree
equalities force $D=\Q(\theta)$, proving (iii).

If $n=\ord(\theta)$, then
\[
 \ord(\theta^r)=\frac{n}{(n,r)}.
\]
Writing $a=(n,r)$ proves \eqref{eq:admissible-orders}, and the converse
is immediate.  The product in \eqref{eq:order-block} is Galois
invariant.  A root $t_0$ of one of its factors has, by (iii), degree
\[
 [\Q(t_0):\Q]=N\varphi(ka),
\]
which is the degree of $F_{k,a}$.  Hence $F_{k,a}$ is the minimal
polynomial of $t_0$.  Distinct $a$ give disjoint root sets, completing
the proof.
\end{proof}

The arithmetic count in the preceding theorem is classical; compare
the standard irreducibility criteria for cyclotomic compositions in
\cite{Damianou2011}.  It is included to make the later dynamical
factorization exact.

\begin{corollary}
\label{cor:abstract-factor-count}
For all sufficiently large $k$, the number of irreducible factors of
$\Psi_{r,k}$ over $\Q$ is
\[
 \mathfrak f_r(k)=
 \prod_{\substack{p^\alpha\parallel r\\p\nmid k}}(\alpha+1).
 \numberthis\label{eq:factor-count}
\]
In particular,
\[
 \Psi_{r,k}\text{ is irreducible over }\Q
 \quad\Longleftrightarrow\quad
 \operatorname{rad}(r)\mid k.
 \numberthis\label{eq:radical-criterion}
\]
Consequently the set of such $k$ has natural density
$1/\operatorname{rad}(r)$.
\end{corollary}

\begin{proof}
For a prime power $p^\alpha\parallel r$, condition
$(k,r/a)=1$ forces the exponent of $p$ in $a$ to equal $\alpha$ when
$p\mid k$, while all $\alpha+1$ exponents are allowed when $p\nmid k$.
This proves \eqref{eq:factor-count} and \eqref{eq:radical-criterion}.
The density statement is unchanged by the finite exceptional set.
\end{proof}

We shall also need the same result for a cover that is itself defined
only over a number field.  The finite exceptional set may depend on
that field.

\begin{theorem}[Number-field Kummer--torsion specialization]
\label{thm:kummer-torsion-number-field}
Let $B$ be a number field, let $Z/B$ be a geometrically integral normal
curve, and let $u,v\in B(Z)$ satisfy
\[
 B(Z)=B(u,v),\qquad N=[B(Z):B(v)].
\]
Let $A_B(T,Y)\in B(Y)[T]$ be the monic minimal polynomial of $u$ over
$B(Y)$.  Assume that $v:Z\dashrightarrow\mathbf G_m$ satisfies
\emph{(PB)} and that $u$ separates the torsion fibers of $v$ away from
a finite Galois-stable set.  Fix $r\ge1$, put
\[
 S_r(k)=\{\theta\in\mu_\infty:\ord(\theta^r)=k\},\qquad
 \Psi^B_{r,k}(T)=\prod_{\theta\in S_r(k)}A_B(T,\theta),
\]
whenever all specializations are defined.  Then the following hold for
every sufficiently large $k$.

\begin{enumerate}[label=\textup{(\roman*)}]
\item Over $B^c=B(\mu_\infty)$, the distinct irreducible factors of
      $\Psi^B_{r,k}$ are precisely the $A_B(T,\theta)$, each of degree
      $N$.
\item Over $B$, the irreducible factors are indexed by the
      $\Gal(B^c/B)$-orbits $\mathcal O$ on $S_r(k)$.  The orbit
      $\mathcal O$ gives a factor of degree $N|\mathcal O|$.
\item If $u_0$ is a root of $A_B(T,\theta)$, then
      \[
       \theta\in B(u_0),\qquad
       B(u_0)\cap B^c=B(\theta).
       \numberthis\label{eq:B-intersection}
      \]
\end{enumerate}
For an admissible order $n=ka$, where $a\mid r$ and $(k,r/a)=1$, put
$K_n=\Q(\zeta_n)$ and $D_n=B\cap K_n$.  The order-$n$ block has
$[D_n:\Q]$ irreducible factors over $B$, each of degree
$N[K_n:D_n]$.
\end{theorem}

\begin{proof}
Apply Zannier's theorem \cite[Theorem 2.1]{Zannier2010} over $B$.  In
dimension one its exceptional union of proper torsion cosets is a
finite set of torsion points.  Enlarge it by the finite sets on which
specialization or separation fails.  Thus, for every remaining torsion
point $\theta$, the fiber $v^{-1}(\theta)$ is irreducible over $B^c$
and has degree $N$.  Every $\theta\in S_r(k)$ has order at least $k$,
so all labels are good once $k$ is sufficiently large; separation
makes the resulting factors pairwise distinct.  This proves (i).

Let $P\in v^{-1}(\theta)$ and put $u_0=u(P)$.  If
$\sigma\in\Gal(\overline\Q/B(u_0))$, then
\[
 u(\sigma P)=u_0,\qquad v(\sigma P)=\sigma\theta.
\]
Separation gives $\sigma P=P$, hence $\sigma\theta=\theta$ and
$\theta\in B(u_0)$.  Irreducibility over $B^c$ and over the subfield
$B(\theta)$ gives
\[
 [B^c(u_0):B^c]=N=[B(u_0):B(\theta)].
\]
For $D=B(u_0)\cap B^c$, the compositum degree formula gives
$[B(u_0):D]=N$.  Since $B(\theta)\subseteq D$, the two degree
equalities force $D=B(\theta)$, proving (iii).

For a $\Gal(B^c/B)$-orbit $\mathcal O$, the product of the factors
labelled by $\mathcal O$ belongs to $B[T]$.  If $u_0$ lies above
$\theta\in\mathcal O$, then (iii) gives
\[
 [B(u_0):B]
 =N[B(\theta):B]
 =N|\mathcal O|,
\]
which is the degree of that product.  It is therefore the minimal
polynomial of $u_0$ over $B$.  This proves (ii).  Finally, restriction
identifies the action on primitive $n$-th roots with
$\Gal(K_n/D_n)$.  This action is free, so it has $[D_n:\Q]$ orbits,
each of size $[K_n:D_n]$, proving the last assertion.
\end{proof}
